\section{Boltz-2 architecture: a high-level view}

Boltz-2 combines two related tasks: it first predicts a full-atom
biomolecular complex and then uses the predicted interface to estimate whether
the ligand binds and how strong that binding may be \citep{passaro2025boltz2}.
The model is easiest to understand as four connected components: a trunk, a
diffusion decoder, a confidence module, and an affinity module.

\begin{center}
\fbox{\parbox{0.92\textwidth}{\centering
Molecular inputs
$\longrightarrow$ representation trunk
$\longrightarrow$ full-atom diffusion
$\longrightarrow$ candidate 3D complexes
$\longrightarrow$ confidence ranking
$\longrightarrow$ binder probability + continuous affinity score}}
\end{center}

\subsection{Inputs and molecular representation}

For a protein--ligand example, the basic inputs are the protein sequence and a
ligand description such as SMILES. The model can also use a protein
multiple-sequence alignment, structural templates, experimental-method labels,
and user-provided contact or pocket constraints. An atom-attention encoder
combines detailed atomic features into token representations. Protein residues
are represented as residue-level tokens, whereas ligand heavy atoms are
represented individually.

The trunk maintains two complementary learned objects:

\begin{itemize}
  \item a \emph{single representation} $s_i$ describing each token; and
  \item a \emph{pair representation} $z_{ij}$ describing the relationship
        between every pair of tokens.
\end{itemize}

Template and MSA information are incorporated before a 64-layer PairFormer
updates these representations. Recycling sends the updated representations
through the trunk repeatedly so that sequence, molecular identity, and
putative geometry can be refined together.

\subsection{Full-atom diffusion decoder}

The trunk does not itself provide the final coordinates. Boltz-2 begins the
structure-generation stage with noisy, effectively random atomic coordinates
and applies a learned denoising module repeatedly. At a schematic level,
\[
  \mathbf{x}_{t}
  \longrightarrow
  \mathbf{x}_{t-1}
  \longrightarrow \cdots \longrightarrow
  \mathbf{x}_{0},
\]
where the trunk representations condition every denoising step. The final
$\mathbf{x}_0$ contains explicit coordinates for the modeled atoms in the
protein--ligand complex. Because the initial coordinates and sampling process
are stochastic, several candidate complexes can be generated from the same
input.

Optional steering potentials can bias this reverse-diffusion process toward
user constraints or improved physical plausibility. These potentials guide
sampling; they do not turn the neural decoder into a molecular-dynamics or
free-energy calculation.

\subsection{Confidence module}

A separate, smaller PairFormer examines the final trunk representation and a
predicted coordinate set. It estimates confidence quantities such as local and
pairwise errors and interface confidence. These values help rank alternative
structure samples. Confidence answers ``how trustworthy is this predicted
geometry?''; it is not itself a measurement of binding strength.

\subsection{Affinity module}

The affinity calculation uses both the learned trunk representation and the
predicted structure. In the reported inference protocol, five structures are
generated and the candidate with the highest protein--ligand interface
confidence is selected. A specialized PairFormer then focuses on
protein--ligand and intra-ligand pairs, combines $z_{ij}$ with a distogram
derived from the selected coordinates, and pools the interface into one global
representation.

Two multilayer-perceptron heads produce distinct outputs:

\begin{enumerate}
  \item a binding-likelihood probability intended for binder/non-binder
        classification; and
  \item a continuous, logarithmic affinity value intended mainly for ranking
        related ligands against the same target.
\end{enumerate}

The continuous head is trained on a mixture of $K_i$, $K_d$, and IC$_{50}$
measurements standardized to micromolar units. It should be interpreted as an
IC$_{50}$-like measure of binding strength, not as an equilibrium constant from
one uniform assay and not as a calculated $\Delta G_{\mathrm{bind}}$.

\subsection{What Boltz-2 gives us---and what it does not}

Boltz-2 therefore provides two kinds of predictions that must remain separate:

\begin{center}
\begin{tabular}{p{0.27\textwidth}p{0.61\textwidth}}
\toprule
\textbf{Output family} & \textbf{Scientific interpretation} \\
\midrule
Structure and confidence & A sampled full-atom complex and the model's
estimated uncertainty in that geometry \\
Binder likelihood & A learned classification probability for interaction \\
Continuous affinity & A heterogeneous-assay, IC$_{50}$-like ranking score \\
\bottomrule
\end{tabular}
\end{center}

The affinity estimate can fail when the model selects the wrong pocket,
misplaces the ligand or protein conformation, or omits important context. The
released affinity module does not explicitly model contributions from water,
ions, essential cofactors, or all binding partners. Full-atom coordinates make
mechanistic inspection possible, but they do not guarantee physical validity
or thermodynamic accuracy.
